\chapter[The World after Rome]{After Rome, the World Did Not Go Dark}
\section{A Gold Coin without Borders}
Pick up something small, scarcely larger than your thumbnail but surprisingly heavy. It is gold.

This solidus was minted around 500 CE. One side bears an emperor's profile, helmeted with a spear across his shoulder; the other shows Victory holding a cross, surrounded by Latin lettering. Notice three things. First, it was minted not in Rome but in Constantinople, today's Istanbul. Second, according to textbooks, the Roman Empire had already been dead for over twenty years. Third, coins of this reliable weight and purity remained hard currency throughout the known world for centuries, recognised by merchants from Scandinavian hoards to southern Indian ports.

More surprisingly, archaeologists have found Byzantine gold coins and their imitations in Chinese tombs in Xinjiang, Shaanxi, Ningxia, and elsewhere. They travelled the Silk Road with caravans before joining Northern Dynasties or Tang tomb occupants underground.

A fourth point hides in the design. The period's ``barbarian'' kings in Italy, Gaul, and North Africa also minted gold, often bearing not their own faces but Constantinople's emperor, even with his name engraved awkwardly. The Ostrogothic king Theodoric ruled Italy for over thirty years using Roman procedures, employing Roman scholars as ministers, repairing aqueducts, and staging games. Picture these supposed destroyers competing to stamp the destroyed empire's ruler onto their money and govern Roman lands in Roman fashion. What sort of destroyers behave like that?

Hold the coin and remember three contradictions: a dead empire still mints money; its money remains international currency; and that money travels into Chinese graves.

Did it die, then, or not?

\section{476: That Quiet Autumn}
Textbooks give a famous date: 476, the fall of the Western Roman Empire. Almost every school pupil has memorised it. It sounds as though it should come with fire, smoke, collapsing walls, and screams.

Come and look at the scene.

Not Rome, but Ravenna in northeastern Italy, wrapped in marsh mist. The final western court sheltered here behind easily defended wetlands. Morning fog rises over water; the air smells of the sea and woodsmoke. Grain ships enter and leave the docks as usual: residents are used to changing emperors. It is autumn 476. The central figure is a teenager with interminable official titles and a name people mock: Romulus Augustus, echoing Rome's legendary founder Romulus and its first emperor Augustus. They turn Augustus into the diminutive Augustulus, ``little Augustus.'' Even the final emperor's name is a joke.

That autumn, Odoacer, a commander of Germanic mercenaries, tires of fighting on an emperor's behalf, leads troops into Ravenna, and deposes the boy.

Then what? Listen carefully: silence. No massacre, burning palace, or street battle. Odoacer packs up the imperial robes and insignia and sends them to Emperor Zeno in Constantinople, with a message roughly saying that the west needs no emperor and he will administer Italy for Zeno. The boy is neither beheaded nor exiled to an island, but retired to a villa on the Bay of Naples with a pension. Imperial collapse resembles an executive departure: surrender your badge, collect compensation, leave respectably.

Outside Ravenna, marsh mist rises as usual. Tax collectors find a changed name on their accounts, soldiers new designs on their standards, and farmers discover nothing: wheat must still be planted and rent paid.

If this is Rome's fall, something already seems amiss. How did an event barely felt by participants become humanity's most famous collapse?

\section{Who Said Rome Had Fallen?}
State the conflict before reaching conclusions.

For something to count as a fall, something must stop. After 476, what stopped?

Emperors did, but only in the west. Constantinople's emperors continued for almost another millennium, until 1453. Eastern inhabitants never thought theirs another country. They called themselves Romans, Rhomaioi in Greek, and their state the Roman Empire until its final day. ``Byzantine Empire'' is a later European scholarly name they never used themselves.

Did law stop? Decades later, Justinian, reigning from 527 to 565, commissioned the Corpus Juris Civilis, assembling nearly a millennium of Roman law into a structure that later underpinned continental European legal systems. Western barbarian kingdoms retained Roman law too: the Visigoths issued a compilation for their Roman subjects in 506; in Frankish courts, Romans and Franks litigated under their respective rules. Law changed masters, not ceased.

Did language stop? Latin continued in western kitchens, markets, and beside cradles, gradually changing accents into French, Spanish, Italian, and Portuguese. Today's English also contains many Latin words inherited from that supposed death.

Did cities stop? Many genuinely declined. Rome fell from perhaps over a million residents at its peak to tens of thousands in the early Middle Ages. Britain fared worse: towns emptied and even money disappeared for over a century. ``Nothing changed'' is therefore wrong too.

Security quietly changed location. Roman security came from the state: frontier troops, travelling judges, and theoretically someone responsible when things went wrong. After the western court vanished, security became local. Whoever could gather armed men, open grain stores, or provide a castle refuge became the practical government. Manors and local lords dominated western rural power for centuries; peasants surrendered labour and produce for walls that could close during war. This was not a designed system but a way of living assembled from broken order. Meanwhile the church wove another network: bishops remained in declining cities; rural monasteries relieved disasters, taught, sheltered travellers, and mediated disputes, performing many vanished imperial functions. Armed lordship and the church together caught ordinary people falling through political disintegration.

The question is clearer: after 476, \textbf{politics} fragmented, much of \textbf{everyday life} continued, and \textbf{some places} suffered terribly. Which level does ``Rome fell'' describe? More importantly, who packaged these developments as a fall, and why?

\section{One Event, Five Pairs of Ears}
Let contemporaries speak. The following four passages are not quotations from historical records but reconstructions of four positions based on real circumstances. Listen for each speaker's priorities.

\textbf{A Constantinopolitan official}: ``Fall? Absurd. The empire is fine. Its capital is New Rome; its emperor mints gold, collects taxes, revises laws, and builds great churches. The west consists of provinces temporarily occupied by Germans, to be recovered.'' This was not mere bravado. Justinian later sent Belisarius to reconquer North Africa and Italy. Ironically, nearly twenty years of fighting to recover Rome devastated Italy far more than its fall. The homeland paid the bill for its liberation.

\textbf{An old Roman landowner in Gaul}: ``My grandfather paid the emperor; I pay a Frankish king. The tax bill's name changes, not the land or vines. My son must learn some Frankish, and Germanic rules have entered the courts. Times are rougher, certainly, but life remains life.''

\textbf{A Gothic veteran settled on Roman land}: ``We came not to destroy Rome but because we admire it. We wear Roman-style clothes, employ Roman stewards, and put the emperor on our coins. Destroy everything and who pays us rent? Whom do we impress?'' This too reflects reality: barbarian kingdoms competed to inherit Rome more than to bury it.

\textbf{A Silk Road gold trader}: ``I do not understand your fall. This solidus leaves Constantinople, passes through Persian and Sogdian hands, and is recognised in Chang'an. Western roads are harder, with more armed bands and checkpoints, but trade is not dead: it detours, changes hands, and costs more. Ships go where land routes break; caravans go where shipping fails.'' Archaeology supports this: Byzantine gold, Persian silver, and western Asian glass and precious-metal vessels reached Northern Dynasties and Tang tombs.

\textbf{A monk in a scriptorium}: ``The emperor's identity matters little here. I need enough parchment and hands that can hold a pen for a few more years. What you call an ancient classic may survive throughout Europe only in the copy on my desk.''

These five positions make Rome's fall different things: political bankruptcy for a court, a changed landlord for landowners, a takeover for Germans, rerouting for merchants, and quiet preservation largely detached from politics for monks. Remember too that nobody in western Europe during that millennium called their time the Dark Ages. The label came later. From whom? Hold that question.

\section[Thinking Bridge: Checking Era Labels]{Thinking Bridge: Give Period Labels a Check-up}
This tool works today, tomorrow, and whenever you scroll. Historians call it \textbf{periodisation}: dividing continuous time into named ages, golden or dark, orderly or chaotic, old or new. Division is necessary to tell history, but someone applies each label at a particular moment for a purpose. Whenever you meet a grand period label, ask four diagnostic questions:

\begin{enumerate}
\item \textbf{Who applied it?} Contemporaries or descendants centuries later? Dark Age inhabitants never called themselves dark.
\item \textbf{When?} Labellers often want distance from the past. Labels advertise by diminishing predecessors and elevating their own generation.
\item \textbf{For whom, and to what end?} A dark medieval period makes one's own age bright; a chaotic previous dynasty prepares the present dynasty's founding story.
\item \textbf{What does it illuminate or conceal?} ``Dark Ages'' highlights genuinely decaying western monuments while hiding gold coins, watermills, and Constantinople's lights.
\end{enumerate}
Practise on praise too: the check-up does not distinguish compliments from insults. Surely a golden age is real? Ask who labelled it, often nostalgic successors. What does it show? Athenian drama and sculpture, Chang'an's poetry and wine. What does it hide? Athens rested on labour by enslaved people and disenfranchised foreigners; Chang'an's poetry did not reach hungry villages. Gold often plates only a visible minority. This does not deny achievements; it reminds us that flattering labels also edit.

Try Chinese history. After the Han collapsed, division, war, and rival states lasted intermittently for three or four centuries, no less turbulent than post-Roman western Europe. Yet traditional historians did not discard the whole period as dark. They used ``Wei, Jin, and the Northern and Southern Dynasties,'' a dry chronological name. We consequently remember Jian'an poetry, Wang Xizhi's calligraphy, and Northern Wei grottoes: nobody calls those centuries a cultural black hole. Similar histories under different labels produce different memories. Labels do not merely describe history; they edit it.

\section{A Page from 533}
Read a short primary source. In Constantinople in 533, fifty-seven years after Rome's supposed fall, Justinian's commissioned legal textbook, the Institutes, opens with a definition borrowed from the earlier jurist Ulpian:

\begin{quote}
Justice is the constant and enduring will to give each person what is due. (Iustitia est constans et perpetua voluntas ius suum cuique tribuendi.---Justinian's Institutes, Book I.)
\end{quote}
Pause for three seconds. Justice is neither the strong person's interest nor divine will here, but something concrete: \textbf{each receiving their due}. Debts repaid, contracts honoured, damage compensated: not a heavenly slogan but balance in an accountant's room.

The textbook's date explains our title. During the supposedly dark period, in a Roman city deemed not Rome, scholars organised nearly a millennium's scattered laws, precedents, and theories into the Corpus Juris Civilis. Centuries later European scholars rediscovered and examined it, providing a framework for modern civil codes in continental Europe, Latin America, and many other countries. The supposed rupture produced a conduit reaching today.

Look further away. Chinese courts likewise knew that western empire still existed. The Old Book of Tang's account of western peoples describes Fulin, ``also called Daqin,'' the Roman Empire of earlier histories, and records tribute embassies during the Zhenguan reign. Chang'an's historians thus neither knew nor recognised Rome's fall: their map showed the same Daqin with another name and capital, still sending envoys and glassware. One civilisation's death may be absent from its neighbour's history lessons.

Add the opening coin. Byzantine gold continued for centuries after the western fall, maintaining purity and serving Mediterranean and more distant merchants as a settlement standard. Money is history's least flattering witness: nobody uses a dead government's currency as hard money. Coins, laws, and embassies together make a suddenly dark world difficult to defend.

\section[Power Cut or New Lamp?]{Power Cut or Changed Light Bulb? Competing Explanations}
The same facts---no western emperor, a surviving eastern one, shrinking cities, living law, fewer inscriptions, circulating gold---support different explanations. Keep four things distinct: what happened; why, disputed below; contemporaries' understandings, our five listeners; and our evaluation of the Dark Ages label. Here we compare the second category.

\textbf{Explanation one: this was a real, catastrophic collapse.}

Give this case its full strength: it deserves attention because archaeology supports it. Pottery: even ordinary Roman farmers could obtain consistent mass-produced tableware; afterwards many western regions returned to handmade, slow-wheel production. The dramatic quality decline marks worsening lives. Roofs: widespread Roman tiles disappeared in many regions, replaced by thatch, prompting archaeologists to say people had forgotten even roofing. Cities and literacy: Rome shrank to a small fraction of its peak; public writing on stone and walls nearly vanished, literacy retreating into monasteries. Trade: cheap, abundant Mediterranean oil, wine, and fine pottery gave way to a trickle of luxury goods. Britain was most extreme: within a generation or two of legionary withdrawal, towns emptied, minting stopped, brick and tile skills disappeared, and people returned to older roundhouses. Here collapse is an appropriate word. Ordinary people in Britain or rural Gaul could buy fewer, poorer goods, lived in darker houses, and endured harder winters. Calling this collapse is no exaggeration.

Its limitation is measuring western European monuments and pottery, then treating the result as the world. Collapse was real within that range, while Constantinople, Alexandria, Damascus, and later Baghdad and Chang'an shone beyond it. More deeply, it measures everything by Roman material life, with a ruler Rome itself made.

\textbf{Explanation two: not death but prolonged transformation.}

Rome's political skeleton broke, but civilisation exceeded that skeleton. Barbarian kings copied its law; popular Latin became Romance languages; churches inherited administrative networks, calendars, and remnants of schools. Monasteries became civilisation's backup drives. Incidentally, today's AD dating system was devised by a monk calculating Easter in 525, within the supposed darkness. Technology's ledger was mixed too. Aqueduct arches collapsed, but rural watermills multiplied: England's 1086 Domesday Book lists over five thousand, a density never reached in Roman times. Heavy ploughs, horse collars, and horseshoes spread during these centuries, changing what farmers cared about: how deeply they could plough and how many mouths they could feed. On this account, the light did not go out; a dimmer lamp illuminated more homes.

Knowledge preservation belongs in the same ledger. In 800, the Frankish king Charlemagne was crowned emperor in Rome. Over three centuries after disappearing in the west, the title of Roman emperor was retrieved from the museum and worn again: Rome survived as a brand everyone wanted to inherit. More substantially, Charlemagne's court and monasteries undertook mass copying of available ancient books. The oldest surviving complete manuscripts of many Latin classics, including works of Cicero and Tacitus, come from this era's scribes. Without them, classical antiquity might genuinely mean only ruins. Greek works translated and annotated in Arabic followed another preservation route. Copying east and west provided civilisation with double insurance.

This explanation can sound weightless. For generations after 536, transformation was a luxury word. Procopius described that year's sun as dim as during an eclipse; tree rings and ice cores later confirmed a global volcanic dust veil, followed by recurring plague from 541. People facing famine and epidemic needed food, not historical philosophy. Telling disaster's victims that it was really a transition can be offensive.

\textbf{A third voice: perhaps the question itself is wrong.}

A more radical approach replaces ``Did Rome fall?'' with ``Is our idea of one world mistaken?'' From the fifth to tenth centuries, no single Rome-centred world existed whose light depended on Rome. Several regional worlds shone: eastern Rome around Constantinople; Islam around Baghdad, translating Greek philosophy and science into Arabic, through which many Aristotelian works later returned to European classrooms; Indian kingdoms and Southeast Asian maritime networks; and East Asia, where reunification in 589 after three centuries of division led into Sui and Tang flourishing. Europe, even Roman Europe, was only one such world. The question shrinks accordingly: one western European lamp dimmed for centuries, but humanity's room never depended on one lamp.

This also clarifies an event Roman-centred accounts struggle with: the Arabs' sweeping seventh-century rise. Byzantium and Sasanian Persia exhausted one another in warfare from 602 to 628, chiefly across today's Syria and Iraq. When Arab armies moved north, both powers' frontiers and treasuries were overstretched. Seen from Rome alone, the newcomers seem to descend from nowhere; seen across regional worlds, the possibility is almost written on the map. Post-Roman world history was never Rome's sequel alone.

\section[Further Challenge: Inventing the Middle Ages]{Further Challenge: The Middle Ages Were Invented}
Now take on a harder concept. Above the Dark Ages label stands a larger one: the Middle Ages.

Examine ``middle.'' Between what? Two good periods, Greco-Roman antiquity and modernity. In effect, the name means \textbf{a thousand years between those periods, waiting to be endured}. The age's very name apologises.

Who invented it? Not medieval people, none of whom announced that they lived in the Middle Ages, any more than we call ourselves ancient. Fourteenth-century scholars such as the Italian poet Petrarch admired Roman Latin and lamented living after its brilliance faded. Darkness originally described \textbf{their own} period, with nostalgia and inferiority. By the seventeenth century, European scholars divided history into ancient, medieval, and modern, a scheme that entered textbooks worldwide.

Recognising this invention reveals something important: \textbf{periodisation is an argument, not a discovery}. Time flows continuously; people cut the boundary ``antiquity ends in 476.'' Every cut and name implies a stance. Declaring Rome dead in 476 expels Constantinople from Rome; calling the intervening millennium medieval advertises modernity. Even the shining Renaissance is a retrospective label: painters did not sign canvases ``Renaissance work.'' Ages are lived first, named later, then arranged into stories successors need. Eighteenth-century British historian Edward Gibbon's Decline and Fall of the Roman Empire largely blamed barbarians and Christianity, fitting his Enlightenment commitments. His history was also contemporary political argument. Almost every history of decline doubles as its author's manifesto for their own age.

Beware a further misreading. Learning that labels are constructed can suggest that periods are all equivalent and decline imaginary, with every warning merely narrative manipulation. That goes too far. The labels are constructed; the pottery is not. Excavated bowls and roofs, not invention, show poorer western European material life around 500 than around 200. A mature approach requires both hands: one dismantles labels to expose editing, the other respects evidence of genuinely hard times in particular places. ``Everything is narrative'' and ``textbook labels are truth'' are equally lazy.

Look into a nearby mirror. Traditional Chinese histories have their own editing habits: each dynasty begins fresh, ends amid public and heavenly anger, and passes through cycles of order and chaos. This framework trims history into standard founding, consolidation, and collapse dramas, reducing late-dynastic lives to wicked rulers, treacherous ministers, and suffering populations. Western darkness and Chinese end-of-dynasty narratives are two products of the same editing technique. Check period labels vigilantly in both systems.

\section{Everyone Awaits Rome's Collapse}
This old question is busy today. You may have encountered it without recognising it.

News and forums teem with Rome's ghost. Every few years a bestseller argues that a great power repeats its decline; online voices announce an age ending; a struggling industry immediately enters its dark age. The structure is identical: proclaim an age dead, then imply privileged insight into historical laws. Use your check-up: who speaks, what do they want, what appears or disappears under the label? Often decline efficiently packages both anxiety and its supposedly definitive remedy.

A closer echo sits in textbooks still saying that Rome fell in 476 and Europe entered the dark Middle Ages. If you memorised that, you now know its complexity exceeds a page. This does not mean textbooks deceive; every page edits, and a slip of the editor's hand reshapes millions of pupils' memories.

Turn the lens onto yourself. Every few years a generation receives a collective label: lost, sheltered, lying flat. The same machine operates: outsiders bundle tens of millions of different people under a disparaging phrase. Who labels them? Often the older generation holding the microphone. When? Often when it feels control slipping. Why? Sometimes disappointed concern; sometimes evasion, blaming deficient youth for a harder world. What disappears? Internal diversity and the structures causing difficulty. Next time a grand word summarises your generation, examine it as you would the Dark Ages.

Listen for Rome in everyday language too. ``All roads lead to Rome'' assumes Rome as the world's centre and destination, yet our post-fifth-century world had no agreed centre. Language outlasts history books: labels die while their grammar survives. Hearing old worldviews hidden in idioms and headlines is another portable tool.

Finally, archaeology and climate science are still rewriting this history. Tree rings, ice cores, and lake sediments testify for people without written records, revealing the 536 dust veil, epidemic patterns, and harvests. These supposedly undocumented dark centuries are among today's fastest-growing fields of evidence. Half their darkness means only that nobody recorded things. Absence of records does not mean nothing happened, still less that nobody lived seriously there.

\section{Your Question: Monuments or Meals?}
Return to the gold coin.

Minted after Rome's supposed fall, circulating through darkness, and finally buried in China far from Rome, its life contradicts a suddenly dark world.

But this chapter offers no simple reassurance that nothing changed. Many ordinary western Europeans genuinely suffered worsening lives: pottery and roofs do not lie. Yet universities, spectacles, mechanical clocks, AD dating, and roots of many legal ideas emerged from those same dark centuries. Both accounts are real and belong to the same period.

Now the question comes to you, without a standard answer:

Weigh what you hold. Collapse theorists offer pottery and roofs, real evidence of hardship rather than prejudice. Transformation theorists offer gold, law, and watermills, genuinely flowing channels of civilisation rather than whitewash. The third voice offers a map: Europe was one among several regional worlds, and its dimming did not darken humanity's room. You also hold a check-up sheet asking who applied each period label, when, for whom, and what it conceals.

\textbf{What scales should we use to judge an age?} Monuments: city size, temple height, documentary abundance? Or meals: what ordinary people ate, whether illness was affordable, whether children survived? By meals, Rome's glory may shine less and medieval darkness seem less black. Yet measuring only meals, how do we count the light kindled by manuscript-copying monks, gold-minting artisans, and scholars compiling law?

Choose your scales and explain why. A hint: your choice may also reveal how you intend to judge your own age.
